problem:  
Let $\Sigma \subset S^n$ be a smooth, closed immersion of fixed genus $g$. For $p>1$ and $\alpha>0$, define the *perturbed Willmore energy*  
\[
\mathcal{W}_{p,\alpha}(\Sigma) = \int_{\Sigma} \big(|H|^{2p} + \alpha |A|^{2p}\big)\, d\mu,
\]  
where $H$ is the mean curvature vector and $A$ the second fundamental form of $\Sigma$ in $S^n$.  

Suppose $\{\Sigma_p\}_{p>1}$ is a family of immersions of fixed genus $g$ that are critical for $\mathcal{W}_{p,\alpha(p)}$, with parameters $\alpha(p)$ satisfying $\alpha(p)\to 0$ as $p\to1^+$, and such that  
\[
\sup_{p>1} \mathcal{W}_{p,\alpha(p)}(\Sigma_p) < +\infty.
\]
Assume further that, up to Möbius transformations of the ambient sphere $S^n$, the sequence $\{\Sigma_p\}$ converges (in the varifold sense) to a limiting immersion $\Sigma_\infty$.

**Question:**  
Does every such limiting surface $\Sigma_\infty$ necessarily correspond to a Willmore critical immersion in $S^n$?  
If so, can one identify geometric or analytic conditions on $\alpha(p)$ and the family $\{\Sigma_p\}$ that ensure that $\Sigma_\infty$ is *not* Möbius-equivalent to any minimal immersion into a totally geodesic $S^m \subset S^n$?  
Equivalently, does the perturbed functional $\mathcal{W}_{p,\alpha(p)}$ produce, via its critical points and the limit $p\to1^+$, new non-minimal Willmore surfaces in $S^n$?

Why is it a "good" problem:  
This formulation refines the question into a mathematically well-posed research problem within modern geometric analysis. It asks whether a natural perturbation of the Willmore energy by higher-$p$ curvature powers yields Γ-convergence in a way that preserves the Willmore condition in the limit, and whether the perturbation might generate new branches of Willmore-critical immersions. The problem is simple to state—essentially asking about the limiting geometry of critical points of $\mathcal{W}_{p,\alpha(p)}$ as $p\to1$—yet touches on deep analytic phenomena: existence and compactness of critical points for nonlinear fourth-order curvature PDEs, varifold convergence, and bifurcation in moduli spaces of Willmore surfaces. It bridges analytic and geometric viewpoints, potentially revealing new existence mechanisms for Willmore immersions beyond known minimal and Clifford families. This mixture of conceptual simplicity, technical difficulty, and potential to unify several geometric-analytic themes makes it a *good* research-level mathematical problem.